
% LTT Document Template for Diploma/Master/Bachelor Thesis/PhD Thesis
% -----------------------------------------------------
%
% original template made by: Franz Huber [Franz.Huber@ltt.uni-erlangen.de]


\NeedsTeXFormat{LaTeX2e}

%\documentclass[DIV1,liststotoc,ngerman,11pt,a5paper,BCOR5mm]{LTTLayout}
%\documentclass[DIVclassic,ngerman,pointlessnumbers,liststotoc,tocindent,12.7pt,a5paper, nochapterprefix, open=right] {LTTLayout}
%\documentclass[ngerman, 9pt, a5paper, open=right]{LTTLayout}
%\pdfpageheight=210mm %Page size DinA5
%\pdfpagewidth=148mm

\documentclass[ngerman, 11pt, oneside, a4paper, open=right,]{LTTLayout}
\pdfpageheight=297mm %Page size DinA4
\pdfpagewidth=210mm

%\bibliographystyle{abbrvdin} % Literaturverzeichnis nach DIN
%\bibliographystyle{plaindin}
%\bibliographystyle{osajnl_Block}
\bibliographystyle{osajnl_newline}

%\usepackage{german}
\inputencoding{utf8}
\usepackage{hyperref} %Referenzen mit Hyperlinks
\usepackage{tabularx}
\usepackage{hhline}
\usepackage{textcomp}
\usepackage{pbox}
\usepackage{graphicx}
\usepackage{epstopdf}
\usepackage{booktabs}
\DeclareGraphicsExtensions{.pdf}
\setcounter{secnumdepth}{4}
\setcounter{tocdepth}{4}
\usepackage[version=4]{mhchem}
\usepackage{cite}


%\usepackage[comma,authoryear]{natbib}
%\bibliographystyle{natdin}

% Weitere Optionen
% ----------------
%liststotoc --> Abbildungsverzeichnis
%tocleft --> Linksbündiges Inhaltsverzeichnis
%\addtocounter{tocdepth}{-1}


% schnelles Erstellen: nur diese Kapitel aufnehmen
% ------------------------------------------------

%\includeonly{Einleitung}


\begin{document}

\frontmatter
\singlespacing


% Erstellen der Titelseite
% ------------------------

% Title
\thesistitle{Numerische Untersuchung von Flashback-Phänomenen in Wasserstoffflammen}

% Name der Fakultät (nur für Doktorarbeit) 
%\thesisfaculty{Diplomarbeit im Fach ...} % auf Grammatik achten!
\thesisfaculty{Bachelorarbeit im Fach Chemie- und Bioingenieurwesen}

% vorgelegt von
\thesisauthor{Haidar Makke}

% aus (nur für Doktorarbeit)
\thesisbirthplace{Musterstadt}

% Fach (nur für Diplomarbeit)
\thesisdiscipline{Chemie- und Bioingenieurwesen}

% Vorsitzender der Prüfungskommission
% (für Diplomarbeit: Prüfer)
\thesischairman{Prof.~Dr.-Ing. Vorsitzender}

% Erstgutachter
% (für Diplomarbeit: Betreuer 1)
\thesisfirstreferee{Prof. Dr.-Ing. Gutachter 1}

% Zweitgutachter
% (für Diplomarbeit: Betreuer 2)
\thesissecondreferee{Prof. Dr.-Ing. Gutachter 2}

% Datum der Prüfung
% (für Diplomarbeit: Abgabetermin)
\thesisexamdate{Prüfungsdatum}

%\makedisstitlepage % Titelseite für Doktorarbeit
\thesisworknumber{XX-YY}
\thesisimnumber{222111444}
\makedipltitlepage	% Titelseite für Diplomarbeit

% Verzeichnisse
% -------------

\singlespacing
\tableofcontents
\listoffigures
\listoftables
\chapter{Formelzeichen und Abkürzungen}\label{chapter:symbols}


\subsubsection{Formelzeichen -- Lateinische Buchstaben}

\begin{tabularx}{\textwidth}{cXc}
Variable & Bezeichnung & Dimension \\ \\
  $a$              & Temperaturleitfähigkeit                & m$^2$/s \\
  $A$              & Flammenoberfläche                      & m$^2$ \\
  $A_i$            & präexponentieller Faktor der Reaktion $i$ & -- \\
  $c_p$            & spezifische isobare Wärmekapazität     & J/(kg\,K) \\
  $D$              & Diffusionskoeffizient                  & m$^2$/s \\
  $E_\mathrm{a}$   & Aktivierungsenergie                    & J/mol \\
  $g_\mathrm{c}$   & kritischer Geschwindigkeitsgradient    & 1/s \\
  $h$              & spezifische Enthalpie                  & J/kg \\
  $k_\mathrm{f}$   & Geschwindigkeitskoeffizient (Hinreaktion) & -- \\
  $l_\mathrm{t}$   & integrales turbulentes Längenmaß       & m \\
  $\dot{m}$        & Massenstromdichte                      & kg/(m$^2$\,s) \\
  $p$              & Druck                                  & Pa \\
  $q_i$            & Netto-Reaktionsrate der Reaktion $i$   & mol/(m$^3$\,s) \\
  $R_\mathrm{m}$   & universelle Gaskonstante               & J/(mol\,K) \\
  $s_\mathrm{L}$   & laminare Flammengeschwindigkeit        & m/s \\
  $s_\mathrm{T}$   & turbulente Flammengeschwindigkeit      & m/s \\
  $T$              & Temperatur                             & K \\
  $T_\mathrm{ad}$  & adiabate Flammentemperatur             & K \\
  $u$              & Strömungsgeschwindigkeit               & m/s \\
  $u^{\prime}$     & turbulente Schwankungsgeschwindigkeit  & m/s \\
  $V_k$            & Diffusionsgeschwindigkeit der Spezies $k$ & m/s \\
  $\dot{V}$        & Volumenstrom                           & m$^3$/s \\
  $W_k$            & molare Masse der Spezies $k$           & kg/mol \\
  $x$              & Ortskoordinate                         & m \\
  $Y_k$            & Massenbruch der Spezies $k$            & -- \\
\end{tabularx}


\subsubsection{Formelzeichen -- Griechische Buchstaben}

\begin{tabularx}{\textwidth}{cXc}
Variable & Bezeichnung & Dimension \\ \\
  $\alpha$              & Temperaturexponent                     & -- \\
  $\beta$               & Druckexponent                          & -- \\
  $\delta_\mathrm{L}$   & laminare Flammendicke                  & m \\
  $\delta_\mathrm{q}$   & Löschabstand                           & m \\
  $\eta$                & Kolmogorov-Längenmaß                   & m \\
  $\kappa$              & Streckungsrate                         & 1/s \\
  $\lambda$             & Wärmeleitfähigkeit                     & W/(m\,K) \\
  $\lambda$             & Luftzahl                               & -- \\
  $\mathcal{L}$         & Markstein-Länge                        & m \\
  $\nu^{\prime}$        & stöchiometrischer Koeffizient (Edukte) & -- \\
  $\nu^{\prime\prime}$  & stöchiometrischer Koeffizient (Produkte) & -- \\
  $\rho$                & Dichte                                 & kg/m$^3$ \\
  $\tau_\mathrm{c}$     & chemisches Zeitmaß                     & s \\
  $\tau_\mathrm{t}$     & turbulentes Zeitmaß                    & s \\
  $\varphi$             & Äquivalenzverhältnis                   & -- \\
  $\dot{\omega}_k$      & molare Bildungsrate der Spezies $k$    & mol/(m$^3$\,s) \\
\end{tabularx}


\subsubsection{Dimensionslose Kennzahlen}

\begin{tabularx}{\textwidth}{cXc}
Kennzahl & Bezeichnung & Definition \\ \\
  $\mathrm{Da}$ & Damköhler-Zahl  & $\tau_\mathrm{t}/\tau_\mathrm{c}$ \\
  $\mathrm{Ka}$ & Karlovitz-Zahl  & $\tau_\mathrm{c}/\tau_\eta$ \\
  $\mathrm{Le}$ & Lewis-Zahl      & $a/D$ \\
  $\mathrm{Ma}$ & Markstein-Zahl  & $\mathcal{L}/\delta_\mathrm{L}$ \\
  $\mathrm{Pe}$ & Péclet-Zahl     & $g_\mathrm{c}\delta_\mathrm{L}/s_\mathrm{L}$ \\
  $\mathrm{Re}$ & Reynolds-Zahl   & $u\,d/\nu$ \\
\end{tabularx}


\subsubsection{Abkürzungen}

\begin{tabularx}{\textwidth}{lX}
BLF   & Boundary Layer Flashback (Grenzschichtrückschlag) \\
CFD   & Computational Fluid Dynamics \\
CIVB  & Combustion Induced Vortex Breakdown \\
LTT   & Lehrstuhl für Technische Thermodynamik \\
MFC   & Mass Flow Controller (Massendurchflussregler) \\
PIV   & Particle Image Velocimetry \\
PLIF  & Planar Laser Induced Fluorescence \\
\end{tabularx}


\subsubsection{Indizes}

\begin{tabularx}{\textwidth}{cX}
  $\mathrm{ad}$ & adiabat \\
  $\mathrm{b}$  & verbrannt (burnt) \\
  $\mathrm{c}$  & kritisch bzw.\ chemisch \\
  $k$           & Laufindex der Spezies \\
  $\mathrm{L}$  & laminar \\
  $\mathrm{st}$ & stöchiometrisch \\
  $\mathrm{T}$  & turbulent \\
  $\mathrm{u}$  & unverbrannt (unburnt) \\
  $0$           & Referenzzustand \\
\end{tabularx}


% Kapitel
% -------

\mainmatter
\setstretch{1.3}

\renewcommand*{\chapterheadstartvskip}{\vspace*{2.4\baselineskip}}% Abstand einstellen
\chapter{Einleitung}\label{chapter:einleitung}

\section{Motivation}\label{sec:motivation}

Die Dekarbonisierung der Energieversorgung zählt zu den zentralen technischen
Herausforderungen der kommenden Jahrzehnte. Während sich die Stromerzeugung aus
erneuerbaren Quellen bereits weitgehend elektrifizieren lässt, bleiben
Hochtemperaturprozesse in der Industrie, die Bereitstellung von Regelleistung
sowie Teile des Schwerlast- und Luftverkehrs auf chemisch gespeicherte Energie
angewiesen. Wasserstoff bietet hier die Möglichkeit, elektrische Energie
speicherbar zu machen und in einer Verbrennungsreaktion wieder freizusetzen,
ohne dass dabei Kohlendioxid entsteht.

Die technische Umsetzung dieser Idee ist jedoch nicht dadurch abgeschlossen,
dass bestehende Brenner mit einem anderen Brennstoff versorgt werden.
Wasserstoff unterscheidet sich in seinen verbrennungstechnischen Eigenschaften
deutlich von den bislang eingesetzten Kohlenwasserstoffen. Die laminare
Flammengeschwindigkeit von stöchiometrischen Wasserstoff-Luft-Gemischen liegt
mit rund 2,9\,m/s etwa um den Faktor sieben über derjenigen von
Methan-Luft-Gemischen \cite{Law_2006, Verhelst_2009}. Hinzu kommen ein sehr
weiter Zündbereich von etwa 4\,Vol.-\% bis 75\,Vol.-\%,
eine um eine Größenordnung geringere Mindestzündenergie und ein Löschabstand,
der weniger als ein Drittel des Wertes von Methan beträgt
\cite{Verhelst_2009, Warnatz_2001}.

Diese Eigenschaften wirken sich unmittelbar auf die Betriebssicherheit
vorgemischter Brenner aus. Die Vormischung von Brennstoff und Oxidator vor der
Reaktionszone ist die technisch etablierte Maßnahme zur Absenkung der
Stickoxidemissionen, da sie lokale Temperaturspitzen im Bereich
stöchiometrischer Verbrennung vermeidet \cite{Lieuwen_2008}. Sie setzt
allerdings voraus, dass die Flamme stabil stromab des Brenneraustritts
verankert bleibt. Genau diese Voraussetzung ist bei Wasserstoff nur noch in
einem deutlich eingeschränkten Betriebsbereich erfüllt.

\section{Problemstellung}\label{sec:problemstellung}

Als Flammenrückschlag oder \emph{Flashback} wird das stromaufwärts gerichtete
Eindringen der Flamme aus dem Brennraum in die Vormischstrecke bezeichnet. Der
Vorgang tritt auf, sobald die Ausbreitungsgeschwindigkeit der Flamme lokal die
entgegengerichtete Strömungsgeschwindigkeit des Frischgases übersteigt. Da die
Vormischstrecke thermisch nicht für den Flammenbetrieb ausgelegt ist, führt ein
Flashback innerhalb kurzer Zeit zu Bauteilschäden und erfordert eine
Schnellabschaltung der Anlage \cite{Lieuwen_2008, Kalantari_2017}.

In der Literatur werden vier grundlegende Mechanismen unterschieden, über die
ein Flammenrückschlag ausgelöst werden kann: der Rückschlag in der
wandnahen Grenzschicht, der Rückschlag durch verbrennungsinduziertes
Wirbelaufplatzen (\emph{combustion induced vortex breakdown}, CIVB), der
Rückschlag in der Kernströmung infolge turbulenter Flammenausbreitung sowie der
akustisch beziehungsweise durch Verbrennungsinstabilitäten getriebene
Rückschlag \cite{Kalantari_2017, Kroner_2003}. Welcher dieser Mechanismen
dominiert, hängt von der Brennergeometrie, dem Drallgrad, dem Druckniveau und
insbesondere vom Äquivalenzverhältnis ab.

Allen Mechanismen gemeinsam ist, dass die laminare Flammengeschwindigkeit
$s_\mathrm{L}$ als charakteristische Geschwindigkeitsskala in die Beschreibung
eingeht. Sie bestimmt über den kritischen Geschwindigkeitsgradienten die
Rückschlagsgrenze in der Grenzschicht, sie skaliert die turbulente
Flammengeschwindigkeit in der Kernströmung und sie geht über die Flammendicke
in die Beschreibung der Flammenstreckung ein. Eine belastbare Vorhersage von
Flashback-Grenzen setzt daher voraus, dass $s_\mathrm{L}$ und die zugehörige
Flammenstruktur für den betrachteten Betriebsbereich hinreichend genau bekannt
sind.

Erschwerend kommt hinzu, dass sich Wasserstoffflammen wegen ihrer niedrigen
Lewis-Zahl grundsätzlich anders verhalten als Kohlenwasserstoffflammen. Magere
Wasserstoff-Luft-Flammen sind thermodiffusiv instabil; sie bilden zellulare
Strukturen aus, deren effektive Ausbreitungsgeschwindigkeit die eindimensional
berechnete laminare Flammengeschwindigkeit deutlich übersteigen kann
\cite{Matalon_2007, Berger_2022, Frouzakis_2015}. Eine rein eindimensionale
Betrachtung liefert damit zwar die notwendige chemische Grundlage, sie stellt
aber gleichzeitig eine konservative untere Schranke dar, deren Gültigkeitsbereich
bekannt sein muss.

\section{Stand der Forschung}\label{sec:stand_der_forschung}

Der Grenzschichtrückschlag wurde erstmals von Lewis und von Elbe über das
Konzept des kritischen Geschwindigkeitsgradienten beschrieben und ist seither
sowohl experimentell als auch numerisch umfassend untersucht worden. Einen
aktuellen Überblick über Mechanismen, Modellansätze und experimentelle Befunde
geben Kalantari und McDonell \cite{Kalantari_2017}. Die Arbeiten von Eichler und
Sattelmayer \cite{Eichler_2012} sowie Baumgartner et al.\ \cite{Baumgartner_2016}
zeigen mittels hochauflösender Geschwindigkeitsmessungen, dass die
rückschlagende Flamme die Grenzschicht nicht passiv durchdringt, sondern über
die Wärmefreisetzung selbst eine Strömungsablösung induziert und sich damit den
Weg stromauf bahnt. Das klassische Modell des kritischen
Geschwindigkeitsgradienten unterschätzt die Rückschlagsneigung deshalb
systematisch.

Numerisch hat sich die Grobstruktursimulation (LES) als tragfähiger Zugang
erwiesen. Endres und Sattelmayer \cite{Endres_2018} zeigen für eingeschlossene
Wasserstoff-Luft-Flammen, dass eine LES mit detaillierter Reaktionskinetik und
differentieller Diffusion die experimentell bestimmten Rückschlagsgrenzen
quantitativ wiedergibt. Ihr zentrales Ergebnis ist zugleich ein
Kriterium: Der Rückschlag setzt ein, sobald die Ausdehnung der lokalen
Ablösegebiete den Löschabstand der Flamme deutlich überschreitet. Fruzza et al.\
\cite{Fruzza_2025} weisen für perforierte Brenner nach, dass zweidimensionale
Rechnungen die Flammendynamik nur unzureichend erfassen und die
Rückschlagsgeschwindigkeiten entsprechend fehlerhaft vorhersagen. Ali und
Bastiaans \cite{Ali_2026} untersuchen den Mechanismus für ultramagere
Wasserstoff-Luft-Freistrahlflammen.

Am Lehrstuhl für Technische Thermodynamik der FAU selbst wurden additiv
gefertigte Vormischbrenner für Wasserstoff mit optischen Messverfahren
verglichen \cite{Faderl_2026}; ergänzend zeigt die Arbeit von Strickling et al.\
\cite{Strickling_ua}, dass eine radiale Gemischschichtung die Rückschlagsneigung
gezielt herabsetzen kann. Diese Arbeiten bilden den unmittelbaren Kontext der
vorliegenden Untersuchung.

Für drallstabilisierte Brenner haben Kröner et al.\ \cite{Kroner_2003} und
Fritz et al.\ \cite{Fritz_2004} den CIVB-Mechanismus als maßgeblichen
Rückschlagspfad identifiziert und quantitative Grenzen für dessen Auftreten
angegeben. Lieuwen et al.\ \cite{Lieuwen_2008} ordnen diese Ergebnisse in den
Kontext der Brennstoffflexibilität ein und zeigen, wie stark die
Betriebsgrenzen bereits durch moderate Wasserstoffbeimischung eingeengt werden.

Auf Seiten der Reaktionskinetik liegen für das System \ce{H2}/\ce{O2}
mittlerweile mehrere detaillierte, breit validierte Mechanismen vor, unter
anderem die Modelle von Ó Conaire et al.\ \cite{OConaire_2004}, Li et al.\
\cite{Li_2004}, Burke et al.\ \cite{Burke_2012} und Kéromnès et al.\
\cite{Keromnes_2013}. Mit der Softwarebibliothek Cantera \cite{Cantera_2023}
steht ein quelloffenes Werkzeug zur Verfügung, mit dem sich auf Basis dieser
Mechanismen eindimensionale, frei propagierende Vormischflammen mit
vertretbarem Rechenaufwand lösen lassen.

Was in der Literatur vergleichsweise selten geschlossen dargestellt wird, ist
die Verbindung zwischen den so berechneten chemischen Kenngrößen und den an
einem konkreten Prüfstand gemessenen Rückschlagsgrenzen. Genau an dieser Stelle
setzt die vorliegende Arbeit an.

\section{Zielsetzung}\label{sec:zielsetzung}

Ziel dieser Arbeit ist es, das Auftreten von Flammenrückschlägen in
Wasserstoffflammen numerisch zu untersuchen und die berechneten Kenngrößen mit
experimentellen Daten des Prüfstands abzugleichen. Im Einzelnen werden folgende
Fragestellungen bearbeitet:

\begin{itemize}
	\item Wie hängen die laminare Flammengeschwindigkeit, die Flammendicke und
	      die adiabate Flammentemperatur von Wasserstoff-Luft-Gemischen vom
	      Äquivalenzverhältnis, von der Vorwärmtemperatur und vom Druck ab?
	\item Welchen Einfluss hat die Wahl des Reaktionsmechanismus auf die
	      berechneten Kenngrößen, und wie groß ist die daraus resultierende
	      Unsicherheit?
	\item Welche Kombinationen aus Volumenstrom und Äquivalenzverhältnis führen
	      am untersuchten Brenner zu einem stabilen, welche zu einem
	      instabilen Betrieb?
	\item Welche aus der Simulation zugänglichen Größen eignen sich als
	      Indikator für einen bevorstehenden Flammenrückschlag?
	\item Wie gut lassen sich die experimentell bestimmten Rückschlagsgrenzen
	      mit den eindimensionalen Rechnungen abbilden, und welche Effekte
	      erfordern eine weitergehende CFD-Betrachtung?
\end{itemize}

Die Untersuchung erfolgt in einem ersten Schritt mit eindimensionalen
Simulationen frei propagierender Vormischflammen in Cantera. Darauf aufbauend
wird bewertet, welche Parameter für eine spätere dreidimensionale
Strömungssimulation, in der Geometrie und Turbulenz explizit berücksichtigt
werden, von Bedeutung sind.

\section{Aufbau der Arbeit}\label{sec:aufbau}

Kapitel \ref{chapter:grundlagen} stellt die theoretischen Grundlagen zusammen.
Behandelt werden die Eigenschaften von Wasserstoff als Brennstoff, die
Grundbegriffe der vorgemischten Verbrennung, die chemische Reaktionskinetik des
Systems \ce{H2}/\ce{O2}, die laminare Flammengeschwindigkeit sowie der Einfluss
von Flammenstreckung und thermodiffusiver Instabilität. Anschließend werden die
Mechanismen des Flammenrückschlags und die numerische Umsetzung in Cantera
beschrieben.

% TODO: Verweise anpassen, sobald die weiteren Kapitel angelegt sind
% Kapitel \ref{chapter:methodik} beschreibt den Versuchsaufbau und das
% Simulationsmodell. In Kapitel \ref{chapter:ergebnisse} werden die Ergebnisse
% dargestellt und diskutiert. Kapitel \ref{chapter:zusammenfassung} fasst die
% Arbeit zusammen und gibt einen Ausblick.

\chapter{Theoretische Grundlagen}\label{chapter:grundlagen}

Dieses Kapitel stellt die für die Arbeit benötigten Grundlagen zusammen.
Ausgehend von den Stoffeigenschaften des Wasserstoffs werden zunächst die
Grundbegriffe der vorgemischten Verbrennung eingeführt
(Abschnitt \ref{sec:vormischverbrennung}). Es folgen die Reaktionskinetik des
Systems \ce{H2}/\ce{O2} (Abschnitt \ref{sec:kinetik}), die laminare
Flammengeschwindigkeit als zentrale Kenngröße (Abschnitt \ref{sec:laminar})
sowie der Einfluss von Flammenstreckung und thermodiffusiver Instabilität
(Abschnitt \ref{sec:streckung}). Abschnitt \ref{sec:turbulent} behandelt die
turbulente Vormischverbrennung, Abschnitt \ref{sec:flashback} die Mechanismen
des Flammenrückschlags. Abschnitt \ref{sec:cantera} beschreibt schließlich das
numerische Modell.

\section{Wasserstoff als Brennstoff}\label{sec:wasserstoff}

Wasserstoff unterscheidet sich in nahezu allen verbrennungsrelevanten
Eigenschaften deutlich von Methan, dem Hauptbestandteil von Erdgas.
Tabelle \ref{tab:stoffwerte} stellt die wichtigsten Kenngrößen gegenüber.

\begin{table}[htbp]
\begin{center}
\caption[Verbrennungstechnische Eigenschaften von Wasserstoff und Methan]{Verbrennungstechnische Eigenschaften von Wasserstoff und Methan bei einem Druck von 1\,bar und einer Temperatur von 298\,K. Zusammengestellt nach \cite{Law_2006, Warnatz_2001, Verhelst_2009}.}
\label{tab:stoffwerte}
\begin{tabular}{lccc}
\toprule[1.2pt]
\textbf{Eigenschaft} & \textbf{Einheit} & \textbf{\ce{H2}} & \textbf{\ce{CH4}} \\
\midrule[1.2pt]\midrule[1.0pt]
Molare Masse                        & g/mol             & 2,016  & 16,04 \\
Unterer Heizwert (massebezogen)     & MJ/kg             & 120    & 50,0  \\
Unterer Heizwert (volumenbezogen)   & MJ/m$^3$          & 10,8   & 35,9  \\
Stöchiometrischer Luftbedarf        & kg/kg             & 34,3   & 17,2  \\
Zündgrenzen in Luft                 & Vol.-\%           & 4--75  & 5--15 \\
Mindestzündenergie                  & mJ                & 0,017  & 0,28  \\
Löschabstand                        & mm                & 0,64   & 2,03  \\
Diffusionskoeffizient in Luft       & cm$^2$/s          & 0,61   & 0,16  \\
Max.\ lam.\ Flammengeschwindigkeit  & m/s               & ca.\ 2,9 & ca.\ 0,40 \\
zugehöriges $\varphi$               & --                & ca.\ 1,7 & ca.\ 1,1 \\
Adiabate Flammentemperatur ($\varphi=1$) & K            & ca.\ 2390 & ca.\ 2225 \\
\bottomrule
\end{tabular}
\end{center}
\end{table}

Für die Frage des Flammenrückschlags sind vier dieser Eigenschaften von
besonderer Bedeutung. Erstens liegt die maximale laminare
Flammengeschwindigkeit rund siebenmal höher als bei Methan, sodass die Flamme
bereits bei deutlich höheren Strömungsgeschwindigkeiten stromauf laufen kann.
Zweitens erlaubt der weite Zündbereich einen Betrieb bei sehr kleinen
Äquivalenzverhältnissen, in denen Kohlenwasserstoffflammen bereits verlöschen
würden; magere Wasserstoffflammen sind gerade dort jedoch besonders anfällig
für zellulare Instabilitäten. Drittens bedeutet der geringe Löschabstand, dass
die Flamme auch in engen wandnahen Bereichen und in Spalten überleben kann, in
denen eine Methanflamme durch Wärmeverlust an die Wand erlischt. Viertens führt
der hohe Diffusionskoeffizient zu einer Lewis-Zahl deutlich unterhalb von eins,
was die in Abschnitt \ref{sec:streckung} beschriebenen thermodiffusiven
Effekte hervorruft.

Der volumenbezogene Heizwert von Wasserstoff beträgt nur etwa
30\,\% desjenigen von Methan. Bei gleicher thermischer Leistung muss folglich
ein deutlich größerer Brennstoffvolumenstrom bereitgestellt werden, was
unmittelbare Konsequenzen für die Auslegung der Zuleitungen und der
Mischungszone hat.

\section{Grundbegriffe der vorgemischten Verbrennung}\label{sec:vormischverbrennung}

\subsection{Äquivalenzverhältnis}\label{subsec:phi}

Die vollständige Oxidation von Wasserstoff mit Luft lässt sich unter der
üblichen Näherung, dass trockene Luft aus 21\,Vol.-\% Sauerstoff und
79\,Vol.-\% Stickstoff besteht, durch die Bruttoreaktionsgleichung

\begin{equation}
	\label{eq:bruttoreaktion}
	\ce{H2} + \frac{1}{2}\left(\ce{O2} + 3{,}76\,\ce{N2}\right)
	\longrightarrow \ce{H2O} + 1{,}88\,\ce{N2}
\end{equation}

beschreiben. Das Äquivalenzverhältnis $\varphi$ setzt das tatsächlich
vorliegende Brennstoff-Oxidator-Verhältnis ins Verhältnis zum
stöchiometrischen Wert,

\begin{equation}
	\label{eq:phi}
	\varphi = \frac{\left(\dot{m}_\mathrm{B}/\dot{m}_\mathrm{L}\right)}
	               {\left(\dot{m}_\mathrm{B}/\dot{m}_\mathrm{L}\right)_\mathrm{st}}
	        = \frac{1}{\lambda},
\end{equation}

wobei $\dot{m}_\mathrm{B}$ und $\dot{m}_\mathrm{L}$ die Massenströme von
Brennstoff und Luft bezeichnen und $\lambda$ die in der Anlagentechnik
gebräuchliche Luftzahl ist. Für $\varphi < 1$ liegt ein mageres, für
$\varphi > 1$ ein fettes Gemisch vor. Im stöchiometrischen Fall
($\varphi = 1$) beträgt der Wasserstoffanteil im Gemisch 29,6\,Vol.-\%.

Da an dem in dieser Arbeit betrachteten Prüfstand die Volumenströme von
Wasserstoff und Luft über Massendurchflussregler eingestellt werden, ist es
zweckmäßig, $\varphi$ direkt aus den Normvolumenströmen zu bestimmen:

\begin{equation}
	\label{eq:phi_volumen}
	\varphi = \frac{\dot{V}_{\ce{H2}}}{\dot{V}_\mathrm{Luft}}
	          \cdot \frac{1}{\left(\dot{V}_{\ce{H2}}/\dot{V}_\mathrm{Luft}\right)_\mathrm{st}}
	        = \frac{\dot{V}_{\ce{H2}}}{\dot{V}_\mathrm{Luft}} \cdot 2{,}38.
\end{equation}

\subsection{Adiabate Flammentemperatur}\label{subsec:adiabat}

Die adiabate Flammentemperatur $T_\mathrm{ad}$ ergibt sich aus der Bedingung,
dass die Enthalpie des Systems bei isobarer, adiabater Umsetzung erhalten
bleibt:

\begin{equation}
	\label{eq:adiabat}
	h\!\left(T_\mathrm{u}, Y_{k,\mathrm{u}}\right)
	= h\!\left(T_\mathrm{ad}, Y_{k,\mathrm{b}}\right).
\end{equation}

Hierin bezeichnen die Indizes $\mathrm{u}$ und $\mathrm{b}$ das unverbrannte
beziehungsweise verbrannte Gas und $Y_k$ die Massenbrüche der beteiligten
Spezies. Bei hohen Temperaturen ist die Dissoziation der Produkte zu
berücksichtigen, sodass $T_\mathrm{ad}$ nicht aus einer einfachen
Enthalpiebilanz mit vollständiger Umsetzung folgt, sondern über die Bedingung
des chemischen Gleichgewichts bestimmt werden muss \cite{Warnatz_2001}. Das
Maximum von $T_\mathrm{ad}$ liegt für Wasserstoff-Luft-Gemische geringfügig auf
der fetten Seite von $\varphi = 1$.

\section{Chemische Reaktionskinetik}\label{sec:kinetik}

\subsection{Elementarreaktionen und Reaktionsgeschwindigkeit}\label{subsec:elementarreaktionen}

Die Bruttoreaktionsgleichung \eqref{eq:bruttoreaktion} gibt lediglich die
Bilanz zwischen Edukten und Produkten wieder. Der tatsächliche Ablauf der
Oxidation setzt sich aus einer Vielzahl von Elementarreaktionen zwischen
Radikalen und stabilen Spezies zusammen. Für eine Elementarreaktion $i$ der
Form

\begin{equation}
	\label{eq:elementarreaktion}
	\sum_{k=1}^{K} \nu_{ki}^{\prime}\,\mathcal{M}_k
	\;\rightleftharpoons\;
	\sum_{k=1}^{K} \nu_{ki}^{\prime\prime}\,\mathcal{M}_k
\end{equation}

folgt die Netto-Reaktionsrate aus dem Massenwirkungsgesetz zu

\begin{equation}
	\label{eq:reaktionsrate}
	q_i = k_{\mathrm{f},i} \prod_{k=1}^{K} \left[X_k\right]^{\nu_{ki}^{\prime}}
	    - k_{\mathrm{r},i} \prod_{k=1}^{K} \left[X_k\right]^{\nu_{ki}^{\prime\prime}},
\end{equation}

wobei $\left[X_k\right]$ die molare Konzentration der Spezies $k$ bezeichnet.
Die Geschwindigkeitskoeffizienten $k_{\mathrm{f},i}$ werden üblicherweise in
der erweiterten Arrhenius-Form

\begin{equation}
	\label{eq:arrhenius}
	k_{\mathrm{f},i}(T) = A_i\,T^{b_i}\,\exp\!\left(-\frac{E_{\mathrm{a},i}}{R_\mathrm{m} T}\right)
\end{equation}

angegeben. Die Bildungsrate der Spezies $k$ ergibt sich durch Summation über
alle $I$ Reaktionen:

\begin{equation}
	\label{eq:speziesquellterm}
	\dot{\omega}_k = \sum_{i=1}^{I}
	\left(\nu_{ki}^{\prime\prime} - \nu_{ki}^{\prime}\right) q_i.
\end{equation}

\subsection{Reaktionspfade der Wasserstoffoxidation}\label{subsec:reaktionspfade}

Das Reaktionssystem \ce{H2}/\ce{O2} umfasst je nach Mechanismus etwa
acht bis zehn Spezies (\ce{H2}, \ce{O2}, \ce{H2O}, \ce{H}, \ce{O}, \ce{OH},
\ce{HO2}, \ce{H2O2} sowie inerte Stoßpartner) und zwanzig bis dreißig
Elementarreaktionen. Es ist damit das am besten untersuchte Verbrennungssystem
überhaupt und bildet zugleich das Grundgerüst jedes
Kohlenwasserstoffmechanismus \cite{Law_2006}.

Von zentraler Bedeutung ist die Kettenverzweigungsreaktion

\begin{equation}
	\label{eq:kettenverzweigung}
	\ce{H} + \ce{O2} \rightleftharpoons \ce{O} + \ce{OH},
\end{equation}

bei der aus einem Radikal zwei entstehen. Sie besitzt eine vergleichsweise hohe
Aktivierungsenergie und bestimmt deshalb maßgeblich die Temperaturabhängigkeit
der Reaktionsgeschwindigkeit. In Konkurrenz dazu steht die
Dreistoßrekombination

\begin{equation}
	\label{eq:kettenabbruch}
	\ce{H} + \ce{O2} + \ce{M} \rightleftharpoons \ce{HO2} + \ce{M},
\end{equation}

die als Kettenabbruch wirkt, da das entstehende \ce{HO2} unter mageren
Bedingungen vergleichsweise reaktionsträge ist. Da Reaktion
\eqref{eq:kettenabbruch} als Dreistoßreaktion druckabhängig ist, Reaktion
\eqref{eq:kettenverzweigung} dagegen nicht, verschiebt sich das Verhältnis
beider Pfade mit steigendem Druck zugunsten des Kettenabbruchs. Aus der
Gleichgewichtsbedingung beider Raten folgt die zweite Explosionsgrenze des
Systems \ce{H2}/\ce{O2}. Dieser Zusammenhang erklärt unmittelbar, weshalb die
laminare Flammengeschwindigkeit von Wasserstoff mit steigendem Druck abnimmt
\cite{Law_2006, Burke_2012}.

Die weitere Umsetzung erfolgt über die Kettenfortpflanzungsreaktionen

\begin{align}
	\label{eq:kettenfortpflanzung}
	\ce{O} + \ce{H2}  &\rightleftharpoons \ce{H} + \ce{OH}, \\
	\ce{OH} + \ce{H2} &\rightleftharpoons \ce{H} + \ce{H2O},
\end{align}

wobei die zweite Reaktion den Hauptanteil der Wasserbildung trägt.

\subsection{Verfügbare Reaktionsmechanismen}\label{subsec:mechanismen}

Für die vorliegende Arbeit kommen mehrere validierte Mechanismen in Betracht.
Der Mechanismus von Ó Conaire et al.\ \cite{OConaire_2004} ist über einen
weiten Bereich von Druck, Temperatur und Äquivalenzverhältnis validiert und
wird häufig als Referenz herangezogen. Li et al.\ \cite{Li_2004} und
Burke et al.\ \cite{Burke_2012} legen den Schwerpunkt auf erhöhte Drücke,
wobei Burke et al.\ insbesondere die Druckabhängigkeit der Reaktion
\eqref{eq:kettenabbruch} überarbeitet haben. Kéromnès et al.\
\cite{Keromnes_2013} erweitern die Validierungsbasis auf Synthesegasgemische.
Der San-Diego-Mechanismus \cite{SanDiego_Mech} zeichnet sich durch einen
bewusst reduzierten Umfang bei gleichzeitig guter Vorhersagequalität aus.

Da die Unterschiede zwischen diesen Mechanismen in der berechneten laminaren
Flammengeschwindigkeit je nach Betriebspunkt mehrere Prozent betragen können,
wird in dieser Arbeit ein Vergleich mehrerer Mechanismen durchgeführt, um die
resultierende Modellunsicherheit zu quantifizieren.

\section{Die laminare Vormischflamme}\label{sec:laminar}

\subsection{Flammenstruktur und laminare Flammengeschwindigkeit}\label{subsec:sl}

Eine frei propagierende, adiabate und eindimensionale Vormischflamme breitet
sich relativ zum unverbrannten Gas mit der laminaren Flammengeschwindigkeit
$s_\mathrm{L}$ aus. Sie ist eine reine Stoffeigenschaft des Gemisches und hängt
ausschließlich von Zusammensetzung, Vorwärmtemperatur $T_\mathrm{u}$ und Druck
$p$ ab.

Die Flamme lässt sich in eine Vorwärmzone, in der das Frischgas durch
Wärmeleitung aus der Reaktionszone aufgeheizt wird, und eine dünne
Reaktionszone unterteilen, in der die Wärmefreisetzung stattfindet. Aus der
Betrachtung des Gleichgewichts zwischen diffusivem Wärmetransport und
chemischer Umsatzrate folgt die auf Mallard und Le Chatelier zurückgehende
Skalierung

\begin{equation}
	\label{eq:mallard}
	s_\mathrm{L} \sim \sqrt{a \cdot \dot{\omega}},
	\qquad
	a = \frac{\lambda}{\rho\,c_p},
\end{equation}

mit der Temperaturleitfähigkeit $a$ und einer charakteristischen
Reaktionsrate $\dot{\omega}$. Die zugehörige laminare Flammendicke wird in
dieser Arbeit über den maximalen Temperaturgradienten definiert,

\begin{equation}
	\label{eq:flammendicke}
	\delta_\mathrm{L} = \frac{T_\mathrm{b} - T_\mathrm{u}}
	{\left|\mathrm{d}T/\mathrm{d}x\right|_\mathrm{max}},
\end{equation}

da diese Definition unmittelbar aus dem berechneten Temperaturprofil folgt und
im Gegensatz zur diffusiven Abschätzung $\delta_\mathrm{D} = a/s_\mathrm{L}$
keine Annahme über konstante Stoffwerte erfordert \cite{Poinsot_2011}.

\subsection{Einflussgrößen}\label{subsec:einflussgroessen}

Die Abhängigkeit der laminaren Flammengeschwindigkeit vom
Äquivalenzverhältnis weist ein ausgeprägtes Maximum auf. Bei
Kohlenwasserstoffen liegt dieses nahe der Stöchiometrie
($\varphi \approx 1{,}05\ldots1{,}1$), da dort die Flammentemperatur maximal
wird. Bei Wasserstoff-Luft-Gemischen verschiebt sich das Maximum dagegen
deutlich auf die fette Seite zu $\varphi \approx 1{,}7$
\cite{Law_2006, Egolfopoulos_1990, Dahoe_2005}. Ursache hierfür ist die hohe
Diffusivität des Wasserstoffmoleküls: In fetten Gemischen diffundiert
Wasserstoff schneller in die Reaktionszone hinein, als Wärme aus ihr
abgeführt wird, sodass die Umsatzrate über den Punkt maximaler
Flammentemperatur hinaus weiter ansteigt.

Die Abhängigkeit von Vorwärmtemperatur und Druck wird üblicherweise durch den
Potenzansatz

\begin{equation}
	\label{eq:sl_skalierung}
	s_\mathrm{L}(T_\mathrm{u}, p) = s_{\mathrm{L},0}
	\left(\frac{T_\mathrm{u}}{T_0}\right)^{\alpha}
	\left(\frac{p}{p_0}\right)^{\beta}
\end{equation}

beschrieben, wobei $s_{\mathrm{L},0}$ der Wert im Referenzzustand
($T_0 = 298\,$K, $p_0 = 1\,$bar) ist. Für Wasserstoff-Luft-Gemische liegt der
Temperaturexponent typischerweise bei $\alpha \approx 1{,}5\ldots2$, der
Druckexponent ist negativ ($\beta < 0$) \cite{Verhelst_2009}. Die
Druckabhängigkeit folgt unmittelbar aus der in
Abschnitt \ref{subsec:reaktionspfade} beschriebenen Konkurrenz zwischen
Kettenverzweigung und Kettenabbruch.

\section{Flammenstreckung und thermodiffusive Instabilität}\label{sec:streckung}

\subsection{Streckungsrate und Markstein-Zahl}\label{subsec:markstein}

Reale Flammen sind gekrümmt und einem Dehnungsfeld ausgesetzt. Die
Streckungsrate $\kappa$ beschreibt die relative zeitliche Änderung eines
Flammenflächenelements $A$,

\begin{equation}
	\label{eq:streckungsrate}
	\kappa = \frac{1}{A}\frac{\mathrm{d}A}{\mathrm{d}t}.
\end{equation}

Für schwache Streckung besteht ein linearer Zusammenhang zwischen der
gestreckten Flammengeschwindigkeit $s_\mathrm{L}^{\kappa}$ und der
ungestreckten Flammengeschwindigkeit $s_\mathrm{L}^{0}$,

\begin{equation}
	\label{eq:markstein}
	s_\mathrm{L}^{\kappa} = s_\mathrm{L}^{0} - \mathcal{L}\,\kappa
	\qquad\text{bzw.}\qquad
	\frac{s_\mathrm{L}^{\kappa}}{s_\mathrm{L}^{0}}
	= 1 - \mathrm{Ma}\cdot\mathrm{Ka},
\end{equation}

mit der Markstein-Länge $\mathcal{L}$, der Markstein-Zahl
$\mathrm{Ma} = \mathcal{L}/\delta_\mathrm{L}$ und der Karlovitz-Zahl
$\mathrm{Ka} = \kappa\,\delta_\mathrm{L}/s_\mathrm{L}^{0}$
\cite{Poinsot_2011, Matalon_2007}.

\subsection{Lewis-Zahl und zellulare Instabilität}\label{subsec:lewis}

Das Vorzeichen der Markstein-Zahl wird im Wesentlichen durch die Lewis-Zahl

\begin{equation}
	\label{eq:lewis}
	\mathrm{Le} = \frac{a}{D} = \frac{\lambda}{\rho\,c_p\,D}
\end{equation}

bestimmt, also durch das Verhältnis von thermischer zu molekularer
Diffusivität des Mangelreaktanden. Für $\mathrm{Le} > 1$ wirkt die Streckung
stabilisierend, für $\mathrm{Le} < 1$ destabilisierend.

Magere Wasserstoff-Luft-Flammen erreichen Lewis-Zahlen von etwa
$\mathrm{Le} \approx 0{,}3$, da Wasserstoff als Mangelreaktand wesentlich
schneller diffundiert, als Wärme geleitet wird. In einem konvex zum Frischgas
gekrümmten Flammenabschnitt fokussiert die Wasserstoffdiffusion den Brennstoff
lokal, während die Wärme divergent abgeführt wird. Die lokale
Flammentemperatur und damit die lokale Ausbreitungsgeschwindigkeit steigen an,
die Auslenkung verstärkt sich. Es entsteht eine selbstverstärkende, zellulare
Flammenstruktur \cite{Matalon_2007, Berger_2022}.

Diese thermodiffusive Instabilität überlagert sich der hydrodynamischen
Darrieus-Landau-Instabilität, die aus der Dichteänderung über die Flammenfront
resultiert und für alle Vormischflammen wirksam ist. Die praktische Konsequenz
ist erheblich: Durch die Vergrößerung der Flammenoberfläche kann die effektive
Ausbreitungsgeschwindigkeit einer zellularen Wasserstoffflamme die
eindimensional berechnete laminare Flammengeschwindigkeit um ein Vielfaches
übersteigen \cite{Frouzakis_2015, Berger_2022}. Eindimensionale Rechnungen
liefern in diesem Bereich also eine untere Schranke für die Rückschlagsneigung.
Dieser Umstand ist bei der Interpretation der Ergebnisse zwingend zu
berücksichtigen.
% TODO: Verweis auf das Ergebniskapitel ergaenzen, sobald dieses angelegt ist

\section{Turbulente Vormischverbrennung}\label{sec:turbulent}

Am untersuchten Brenner liegen turbulente Strömungsbedingungen vor. Die
turbulente Flammengeschwindigkeit $s_\mathrm{T}$ übersteigt $s_\mathrm{L}$,
weil die Turbulenz die Flammenfront faltet und dadurch deren Oberfläche
vergrößert. Nach Damköhler \cite{Damkoehler_1940} gilt im Bereich großskaliger
Turbulenz

\begin{equation}
	\label{eq:damkoehler}
	\frac{s_\mathrm{T}}{s_\mathrm{L}} = \frac{A_\mathrm{T}}{A_\mathrm{L}}
	\approx 1 + C\,\frac{u^{\prime}}{s_\mathrm{L}},
\end{equation}

mit der turbulenten Schwankungsgeschwindigkeit $u^{\prime}$ und einer
empirischen Konstanten $C$.

Zur Einordnung des Verbrennungsregimes dienen die Damköhler-Zahl und die
Karlovitz-Zahl,

\begin{equation}
	\label{eq:da_ka}
	\mathrm{Da} = \frac{\tau_\mathrm{t}}{\tau_\mathrm{c}}
	= \frac{l_\mathrm{t}/u^{\prime}}{\delta_\mathrm{L}/s_\mathrm{L}},
	\qquad
	\mathrm{Ka} = \frac{\tau_\mathrm{c}}{\tau_\eta}
	= \left(\frac{\delta_\mathrm{L}}{\eta}\right)^{2},
\end{equation}

worin $l_\mathrm{t}$ das integrale Längenmaß und $\eta$ das
Kolmogorov-Längenmaß bezeichnen. Im Borghi-Diagramm \cite{Borghi_1985, Peters_2000}
lassen sich damit die Regime der gefalteten Flammenfronten
($\mathrm{Ka} < 1$), der dünnen Reaktionszonen
($1 < \mathrm{Ka} < 100$) und der gerührten Reaktionszonen abgrenzen.

Die Bestimmung von $\delta_\mathrm{L}$ und $s_\mathrm{L}$ aus den in
Abschnitt \ref{sec:cantera} beschriebenen Rechnungen erlaubt es somit, die
Betriebspunkte des Prüfstands im Borghi-Diagramm einzuordnen.

\section{Mechanismen des Flammenrückschlags}\label{sec:flashback}

Ein Flammenrückschlag tritt ein, wenn die Flamme sich lokal schneller stromauf
bewegt, als das Frischgas sie stromab transportiert. Die allgemeine Bedingung
lautet

\begin{equation}
	\label{eq:flashback_bedingung}
	s_\mathrm{T} > u_\mathrm{lokal}.
\end{equation}

Entscheidend ist der Zusatz \emph{lokal}: Der Rückschlag wird nicht durch die
mittlere Strömungsgeschwindigkeit ausgelöst, sondern durch das Vorhandensein
von Bereichen mit reduzierter Geschwindigkeit. Je nachdem, wodurch ein solcher
Bereich entsteht, unterscheidet man die im Folgenden beschriebenen
Mechanismen \cite{Kalantari_2017, Lieuwen_2008}.

\subsection{Rückschlag in der wandnahen Grenzschicht}\label{subsec:blf}

In der wandnahen Grenzschicht sinkt die Strömungsgeschwindigkeit zur Wand hin
auf null ab. Gleichzeitig wird die Flamme durch den Wärmeverlust an die Wand
gelöscht, sobald sie sich der Wand auf weniger als den Löschabstand
$\delta_\mathrm{q}$ nähert. Aus dem Gleichgewicht beider Effekte folgt der von
Lewis und von Elbe eingeführte kritische Geschwindigkeitsgradient

\begin{equation}
	\label{eq:kritischer_gradient}
	g_\mathrm{c} = \left.\frac{\partial u}{\partial y}\right|_\mathrm{Wand,krit}
	\approx \frac{s_\mathrm{L}}{\delta_\mathrm{q}}.
\end{equation}

Ein Rückschlag wird demnach erwartet, sobald der tatsächliche Wandgradient
unter $g_\mathrm{c}$ fällt. Wegen des kleinen Löschabstands von Wasserstoff
(Tabelle \ref{tab:stoffwerte}) liegt $g_\mathrm{c}$ um mehr als eine
Größenordnung über dem Wert von Methan.

Neuere Untersuchungen zeigen allerdings, dass dieses Modell die
Rückschlagsneigung systematisch unterschätzt. Eichler und Sattelmayer
\cite{Eichler_2012} sowie Baumgartner et al.\ \cite{Baumgartner_2016} weisen
mittels hochauflösender Geschwindigkeitsmessungen nach, dass die
Wärmefreisetzung der annähernden Flamme über den Dichtesprung einen
positiven Druckgradienten stromauf erzeugt, der die Grenzschicht ablöst. Die
Flamme schafft sich damit selbst das Rückströmgebiet, in dem sie sich
fortbewegt; der Vorgang ist folglich rückgekoppelt und nicht durch das
ungestörte Strömungsfeld allein beschreibbar.

Endres und Sattelmayer \cite{Endres_2018} bestätigen dieses Bild numerisch und
leiten daraus ein handhabbares Kriterium ab: Der Rückschlag setzt ein, sobald
die Ausdehnung der lokalen Ablösegebiete den Löschabstand der Flamme deutlich
überschreitet. Damit tritt an die Stelle des rein kinematischen Vergleichs von
Gleichung \eqref{eq:kritischer_gradient} eine geometrische Bedingung, die
Wärmefreisetzung und Strömungsablösung miteinander verknüpft. Für die
numerische Vorhersage bedeutet dies, dass sowohl die differentielle Diffusion
als auch die Thermodiffusion im Speziestransport berücksichtigt werden
müssen \cite{Endres_2018}.

\subsection{Verbrennungsinduziertes Wirbelaufplatzen}\label{subsec:civb}

In drallstabilisierten Brennern kann die Wärmefreisetzung der Flamme das
Wirbelaufplatzen stromauf verlagern. Die Flamme reduziert über die
Volumenexpansion die axiale Impulsstromdichte im Wirbelkern und schwächt
gleichzeitig durch Baroklinität die azimutale Vortizität. Beides begünstigt
die Ausbildung eines Rückströmgebiets weiter stromauf, in das die Flamme
nachrückt; der Vorgang setzt sich selbstverstärkend fort. Kröner et al.\
\cite{Kroner_2003} und Fritz et al.\ \cite{Fritz_2004} haben diesen als CIVB
bezeichneten Mechanismus systematisch untersucht und quantitative
Rückschlagsgrenzen angegeben.

\subsection{Rückschlag in der Kernströmung}\label{subsec:core}

Übersteigt die turbulente Flammengeschwindigkeit die axiale
Strömungsgeschwindigkeit im Kern der Strömung, so propagiert die Flamme
unabhängig von Wand- und Dralleffekten stromauf. Dieser Mechanismus wird
insbesondere bei hohen Äquivalenzverhältnissen und geringen
Strömungsgeschwindigkeiten relevant. Da $s_\mathrm{T}$ nach
Gleichung \eqref{eq:damkoehler} unmittelbar mit $s_\mathrm{L}$ skaliert, ist
die laminare Flammengeschwindigkeit hier die maßgebliche Eingangsgröße --
weshalb die eindimensionalen Rechnungen dieser Arbeit für genau diesen
Mechanismus die belastbarste Aussage liefern. Ali und Bastiaans
\cite{Ali_2026} untersuchen den Vorgang für ultramagere Freistrahlflammen und
zeigen, dass die thermodiffusiv verstärkte Flammenausbreitung dort auch bei
sehr kleinen Äquivalenzverhältnissen zum Rückschlag führen kann. Fruzza et al.\
\cite{Fruzza_2025} weisen darauf hin, dass eine Reduktion auf zwei Dimensionen
die Flammendynamik verfälscht -- ein Hinweis, der für eine spätere
CFD-Betrachtung unmittelbar relevant ist.

\subsection{Akustisch und instabilitätsgetriebener Rückschlag}\label{subsec:acoustic}

Verbrennungsinstabilitäten koppeln Wärmefreisetzung und Druckschwankungen im
Brennraum. Die dabei auftretenden Geschwindigkeitsschwankungen können die
momentane Anströmgeschwindigkeit periodisch so weit absenken, dass die Flamme
kurzzeitig stromauf läuft. Dieser Mechanismus ist transient und lässt sich mit
stationären Betrachtungen grundsätzlich nicht erfassen.

\subsection{Kenngrößen zur Bewertung}\label{subsec:kenngroessen}

Zur dimensionslosen Beschreibung der Rückschlagsgrenze hat sich die
Flammen-Péclet-Zahl

\begin{equation}
	\label{eq:peclet}
	\mathrm{Pe} = \frac{g_\mathrm{c}\,\delta_\mathrm{L}}{s_\mathrm{L}}
\end{equation}

etabliert, die eine Übertragung zwischen unterschiedlichen Geometrien und
Betriebsbedingungen erlaubt \cite{Kalantari_2017}. Sowohl $g_\mathrm{c}$ als
auch $\delta_\mathrm{L}$ und $s_\mathrm{L}$ sind aus den in dieser Arbeit
durchgeführten Rechnungen beziehungsweise aus der Prüfstandsgeometrie
zugänglich.

\section{Numerische Modellierung mit Cantera}\label{sec:cantera}

\subsection{Modellgleichungen}\label{subsec:modellgleichungen}

Cantera \cite{Cantera_2023} ist eine quelloffene Programmbibliothek zur
Lösung von Problemen der chemischen Kinetik, der Thermodynamik und des
Stofftransports. Für die vorliegende Arbeit wird das Modell der frei
propagierenden, adiabaten eindimensionalen Vormischflamme
(\texttt{FreeFlame}) verwendet.

Im flammenfesten Bezugssystem lauten die stationären Erhaltungsgleichungen
\cite{Poinsot_2011, Cantera_2023}:

\begin{equation}
	\label{eq:kontinuitaet}
	\dot{m} = \rho\,u = \mathrm{const.},
\end{equation}

\begin{equation}
	\label{eq:spezies}
	\dot{m}\,\frac{\mathrm{d}Y_k}{\mathrm{d}x}
	+ \frac{\mathrm{d}}{\mathrm{d}x}\left(\rho\,Y_k\,V_k\right)
	= \dot{\omega}_k\,W_k,
	\qquad k = 1,\ldots,K,
\end{equation}

\begin{equation}
	\label{eq:energie}
	\dot{m}\,c_p\,\frac{\mathrm{d}T}{\mathrm{d}x}
	= \frac{\mathrm{d}}{\mathrm{d}x}\left(\lambda\,\frac{\mathrm{d}T}{\mathrm{d}x}\right)
	- \sum_{k=1}^{K}\rho\,Y_k\,V_k\,c_{p,k}\,\frac{\mathrm{d}T}{\mathrm{d}x}
	- \sum_{k=1}^{K} \dot{\omega}_k\,h_k\,W_k.
\end{equation}

Hierin sind $V_k$ die Diffusionsgeschwindigkeit, $W_k$ die molare Masse und
$h_k$ die spezifische Enthalpie der Spezies $k$. Der Eigenwert des Problems
ist der Massenstrom $\dot{m}$, aus dem die laminare Flammengeschwindigkeit
unmittelbar folgt:

\begin{equation}
	\label{eq:sl_aus_massenstrom}
	s_\mathrm{L} = \frac{\dot{m}}{\rho_\mathrm{u}}.
\end{equation}

\subsection{Transportmodell}\label{subsec:transport}

Cantera stellt für die Berechnung der Diffusionsgeschwindigkeiten mehrere
Modelle bereit. Das Modell \texttt{mixture-averaged} berechnet für jede Spezies
einen effektiven Diffusionskoeffizienten gegenüber der Mischung, das Modell
\texttt{multicomponent} löst dagegen die vollständige Matrix der binären
Diffusionskoeffizienten und erlaubt zusätzlich die Berücksichtigung der
Thermodiffusion (Soret-Effekt).

Gerade für Wasserstoff ist diese Unterscheidung nicht unerheblich: Wegen der
geringen molaren Masse von \ce{H2} und \ce{H} liefert das gemittelte Modell
insbesondere bei mageren Gemischen merklich abweichende Ergebnisse. In dieser
Arbeit wird daher der Einfluss des Transportmodells auf $s_\mathrm{L}$
explizit untersucht.

\subsection{Numerische Umsetzung und Gitterunabhängigkeit}\label{subsec:numerik}

Das Gleichungssystem wird auf einem adaptiven Gitter mit einem gedämpften
Newton-Verfahren gelöst. Die Gitterverfeinerung wird über die Parameter
\texttt{ratio}, \texttt{slope} und \texttt{curve} gesteuert, welche das
maximale Verhältnis benachbarter Zellweiten sowie die maximal zulässige
Änderung von Steigung und Krümmung der Lösungsvariablen zwischen zwei
Gitterpunkten begrenzen.

Da die berechnete Flammengeschwindigkeit von diesen Parametern abhängt, ist
vor der eigentlichen Parameterstudie eine Gitterunabhängigkeitsstudie
durchzuführen. Ebenso ist sicherzustellen, dass das Rechengebiet hinreichend
groß gewählt wird, damit die Randbedingungen die Lösung nicht beeinflussen.
Beide Punkte werden im Methodikkapitel behandelt.
% TODO: Verweis auf das Methodikkapitel ergaenzen, sobald dieses angelegt ist

% TODO: Verweis auf das Methodikkapitel korrigieren, sobald dieses angelegt ist

%\chapter{Physikalische Grundlagen}\label{chapter:physikalische_grundlagen}

Im folgenden Kapitel sollen die Grundlagen erl?utert werden, auf denen diese Vorlage und sp?tere Arbeit basiert.
Zuerst wird jedoch ein einziges Paper zitiert, um auch ein Literaturverzeichnis zu erstellen \cite{Oltmann_2010} and \cite{Oltmann_2012}.

\section{Erstes Unterkapitel}\label{sec:erstes_unterkapitel}

Hier steht das erste Unterkapitel des Kapitels \ref{chapter:physikalische_grundlagen}. Es tr?gt auch hier die Nummer \ref{sec:erstes_unterkapitel}.


\subsection{Erstes Unterunterkapitel}\label{subsec:erstes_unterunterkapitel}
Hier werden zwei Gleichungen dargestellt wobei die Gleichung (\ref{eq:kugelwelle}) die Kugelwelle beschreibt.

\begin{equation}
	\label{eq:kugelwelle}
		u_K(\mathbf{r},t)=\frac{u_0}{r} \cdot e^{ikr-\omega t},
\end{equation}

\begin{equation}
	\label{eq:ebene_welle}
		u_E(\mathbf{r},t)=u_0 \cdot e^{i\mathbf{kr}-\omega t}.
\end{equation}

Jetzt folgt eine EPS-Grafik durch den Befehl (\textbackslash PictureCaptionHeight\{...\}\{...\}\{...\}\{...\}). Dieser Befehl braucht vier ?bergabewerte: Den Dateinahmen, den Titel der Abbildung im Abbildungsverzeichnis, der Titel direkt unterhalb der Abbildung im Text, sowie die Bildh?he in beliebigen Einheiten, hier beispielsweise 35\,mm. Den Befehl gibt es analog f?r die Bildbreite (\textbackslash PictureCaptionWidth\{...\}\{...\}\{...\}\{...\}), wobei hier der vierte Parameter die Bildbreite ist. Der Befehl geht nur f?r EPS-Dateien, die Endung '.eps' im Dateinamen entf?llt. Die Grafik liegt im Bildordner, dieser wird im cls-file aktuel auf  den Ordner 'Images' definiert (\textbackslash def \textbackslash scr@bildverzeichnis\{images\}). Das Label der Grafik ist automatisch der Bildname. Da die Grafik zu gro? ist, springt sie auf die n?chste Seite.

\PictureCaptionHeight{beugung_einfachspalt}{Beugung einer ebenen Welle am Einfachspalt}{Beugung einer senkrecht auf einen Spalt treffenden ebenen Welle und Beispiel f?r ein resultierendes Beugungsbild.}{35mm}




\subsection{Zweites Unterunterkapitel}\label{sucsec:zweites_unterunterkapitel}

Jetzt kommen zwei Grafiken nebeneinander mit gleicher Bildh?he durch den Befehl (\textbackslash TwoPicturesHeight\{Name1\}\{Name2\}\{Titel Abb.Verz.\}\{Bild-unterschrift\}\{Abstand hor.\}\{Bildh?he\}) mit sechs ?bergabewerten: Zwei Bildnamen, die Unterschriften f?r Abbildungsverzeichnis und Bildunterschrift, sowie den horizontalen Abstand und die jeweilige Breite der Bilder.

\TwoPicturesHeight{airy_scheibchen}{rechtecksblende}{Beugungsbild Airy und Rechtecksblende}{Hier sieht man sehr deutlich zwei Bilder nebeneinander und den Unterschied zwischen Blendenformen}{10mm}{35mm}


Jetzt kommen drei Grafiken nebeneinander mit gleicher H?he, der Befehl funktioniert analog, das Bild erscheint aus Platzgr?nden eine Seite weiter.\newpage

\ThreePicturesHeight{beugung_einfachspalt_rotiert}{airy_scheibchen}{rechtecksblende}{Beugungsbild Einfachspalt Rotiert, Airy und Rechtecksblende}{Hier sieht man sehr deutlich drei Bilder nebeneinander und den Unterschied zwischen Blendenformen sowie den gedrehten Einfachspalt}{10mm}{25mm}

Und jetzt noch eine 2x2 Darstellung:

\FourPicturesBoxHeight{rechtecksblende}{airy_scheibchen}{beugung_einfachspalt_rotiert}{beugung_einfachspalt}{Vier Bilder als Block}{Hier sieht man sch?n vier Bilder neben- und untereinander im Block.}{5mm}{25mm}

Analoge Befehle gibt es f?r die jeweiligen Bildbreiten, dazu cls-file durchsehen. Weitere Definitionen sind gerne selbst zu erg?nzen.

Jetzt kommt eine beispielhafte  Tabelle. Hier muss jeder selbst anpassen um die gew?nschte Zeilen und Spaltenanzahl zu erhalten.

\begin{table}
\begin{center}
\caption[Kurzer Titel f?r Tabellenverzeichnis]{Titel f?r Tabelle selbst, welcher nun auch etwas l?nger sein darf. Bei mehreren Zeilen wird es unsch?n.}
\label{tab:tab1}
\begin{tabular}{ccc}
%\hline
\toprule[1.2pt]
\textbf{Bezeichnung 1} & \textbf{Bezeichnung 2} & \textbf{Bezeichnung 3} \\  %\hline
\midrule[1.2pt]\midrule[1.0pt]
Wert 11       & Wert 12     & Wert 13  	    \\
Wert 21       & Wert 22     & Wert 23  	    \\
Wert 31       & Wert 32     & Wert 33  	    \\
\hline
\end{tabular}
\end{center}
\end{table}

Hier kommt jetzt massenhaft Text. Der ist viel. Und kurz. Oder auch l?nger. Und kam Teilweise oben schon mal vor. So wie dieses hier: 'Im folgenden Kapitel sollen die physikalischen Grundlagen erl?utert werden, auf denen die Arbeit basiert'.  Und dann geht es wieder von norne los. Hier kommt jetzt massenhaft Text. Der ist viel. Und kurz. Oder auch l?nger. Und kam Teilweise oben schon mal vor. So wie dieses hier: 'Im folgenden Kapitel sollen die physikalischen Grundlagen erl?utert werden, auf denen die Arbeit basiert'.  Und dann geht es wieder von norne los. Hier kommt jetzt massenhaft Text. Der ist viel. Und kurz. Oder auch l?nger. Und kam Teilweise oben schon mal vor. So wie dieses hier: 'Im folgenden Kapitel sollen die physikalischen Grundlagen erl?utert werden, auf denen die Arbeit basiert'.  Und dann geht es wieder von norne los. Hier kommt jetzt massenhaft Text. Der ist viel. Und kurz. Oder auch l?nger. Und kam Teilweise oben schon mal vor. So wie dieses hier: 'Im folgenden Kapitel sollen die physikalischen Grundlagen erl?utert werden, auf denen die Arbeit basiert'.  Und dann geht es wieder von norne los. Hier kommt jetzt massenhaft Text. Der ist viel. Und kurz. Oder auch l?nger. Und kam Teilweise oben schon mal vor. So wie dieses hier: 'Im folgenden Kapitel sollen die physikalischen Grundlagen erl?utert werden, auf denen die Arbeit basiert'.  Und dann geht es wieder von norne los. Hier kommt jetzt massenhaft Text. Der ist viel. Und kurz. Oder auch l?nger. Und kam Teilweise oben schon mal vor. So wie dieses hier: 'Im folgenden Kapitel sollen die physikalischen Grundlagen erl?utert werden, auf denen die Arbeit basiert'.  Und dann geht es wieder von norne los. Hier kommt jetzt massenhaft Text. Der ist viel. Und kurz. Oder auch l?nger. Und kam Teilweise oben schon mal vor. So wie dieses hier: 'Im folgenden Kapitel sollen die physikalischen Grundlagen erl?utert werden, auf denen die Arbeit basiert'.  Und dann geht es wieder von norne los. Hier kommt jetzt massenhaft Text. Der ist viel. Und kurz. Oder auch l?nger. Und kam Teilweise oben schon mal vor. So wie dieses hier: 'Im folgenden Kapitel sollen die physikalischen Grundlagen erl?utert werden, auf denen die Arbeit basiert'.  Und dann geht es wieder von norne los. Hier kommt jetzt massenhaft Text. Der ist viel. Und kurz. Oder auch l?nger. Und kam Teilweise oben schon mal vor. So wie dieses hier: 'Im folgenden Kapitel sollen die physikalischen Grundlagen erl?utert werden, auf denen die Arbeit basiert'.  Und dann geht es wieder von norne los. Hier kommt jetzt massenhaft Text. Der ist viel. Und kurz. Oder auch l?nger. Und kam Teilweise oben schon mal vor. So wie dieses hier: 'Im folgenden Kapitel sollen die physikalischen Grundlagen erl?utert werden, auf denen die Arbeit basiert'.  Und dann geht es wieder von norne los. Hier kommt jetzt massenhaft Text. Der ist viel. Und kurz. Oder auch l?nger. Und kam Teilweise oben schon mal vor. So wie dieses hier: 'Im folgenden Kapitel sollen die physikalischen Grundlagen erl?utert werden, auf denen die Arbeit basiert'.  Und dann geht es wieder von norne los. Hier kommt jetzt massenhaft Text. Der ist viel. Und kurz. Oder auch l?nger. Und kam Teilweise oben schon mal vor. So wie dieses hier: 'Im folgenden Kapitel sollen die physikalischen Grundlagen erl?utert werden, auf denen die Arbeit basiert'.  Und dann geht es wieder von norne los. Hier kommt jetzt massenhaft Text. Der ist viel. Und kurz. Oder auch l?nger. Und kam Teilweise oben schon mal vor. So wie dieses hier: 'Im folgenden Kapitel sollen die physikalischen Grundlagen erl?utert werden, auf denen die Arbeit basiert'.  Und dann geht es wieder von norne los. Hier kommt jetzt massenhaft Text. Der ist viel. Und kurz. Oder auch l?nger. Und kam Teilweise oben schon mal vor. So wie dieses hier: 'Im folgenden Kapitel sollen die physikalischen Grundlagen erl?utert werden, auf denen die Arbeit basiert'.  Und dann geht es wieder von norne los. Hier kommt jetzt massenhaft Text. Der ist viel. Und kurz. Oder auch l?nger. Und kam Teilweise oben schon mal vor. So wie dieses hier: 'Im folgenden Kapitel sollen die physikalischen Grundlagen erl?utert werden, auf denen die Arbeit basiert'.  Und dann geht es wieder von norne los. 


Hier kommt jetzt massenhaft Text. Der ist viel. Und kurz. Oder auch l?nger. Und kam Teilweise oben schon mal vor. So wie dieses hier: 'Im folgenden Kapitel sollen die physikalischen Grundlagen erl?utert werden, auf denen die Arbeit basiert'.  Und dann geht es wieder von norne los. Hier kommt jetzt massenhaft Text. Der ist viel. Und kurz. Oder auch l?nger. Und kam Teilweise oben schon mal vor. So wie dieses hier: 'Im folgenden Kapitel sollen die physikalischen Grundlagen erl?utert werden, auf denen die Arbeit basiert'.  Und dann geht es wieder von norne los. Hier kommt jetzt massenhaft Text. Der ist viel. Und kurz. Oder auch l?nger. Und kam Teilweise oben schon mal vor. So wie dieses hier: 'Im folgenden Kapitel sollen die physikalischen Grundlagen erl?utert werden, auf denen die Arbeit basiert'.  Und dann geht es wieder von norne los. Hier kommt jetzt massenhaft Text. Der ist viel. Und kurz. Oder auch l?nger. Und kam Teilweise oben schon mal vor. So wie dieses hier: 'Im folgenden Kapitel sollen die physikalischen Grundlagen erl?utert werden, auf denen die Arbeit basiert'.  Und dann geht es wieder von norne los. Hier kommt jetzt massenhaft Text. Der ist viel. Und kurz. Oder auch l?nger. Und kam Teilweise oben schon mal vor. So wie dieses hier: 'Im folgenden Kapitel sollen die physikalischen Grundlagen erl?utert werden, auf denen die Arbeit basiert'.  Und dann geht es wieder von norne los. Hier kommt jetzt massenhaft Text. Der ist viel. Und kurz. Oder auch l?nger. Und kam Teilweise oben schon mal vor. So wie dieses hier: 'Im folgenden Kapitel sollen die physikalischen Grundlagen erl?utert werden, auf denen die Arbeit basiert'.  Und dann geht es wieder von norne los. Hier kommt jetzt massenhaft Text. Der ist viel. Und kurz. Oder auch l?nger. Und kam Teilweise oben schon mal vor. So wie dieses hier: 'Im folgenden Kapitel sollen die physikalischen Grundlagen erl?utert werden, auf denen die Arbeit basiert'.  Und dann geht es wieder von norne los. Hier kommt jetzt massenhaft Text. Der ist viel. Und kurz. Oder auch l?nger. Und kam Teilweise oben schon mal vor. So wie dieses hier: 'Im folgenden Kapitel sollen die physikalischen Grundlagen erl?utert werden, auf denen die Arbeit basiert'.  Und dann geht es wieder von norne los. Hier kommt jetzt massenhaft Text. Der ist viel. Und kurz. Oder auch l?nger. Und kam Teilweise oben schon mal vor. So wie dieses hier: 'Im folgenden Kapitel sollen die physikalischen Grundlagen erl?utert werden, auf denen die Arbeit basiert'.  Und dann geht es wieder von norne los. Hier kommt jetzt massenhaft Text. Der ist viel. Und kurz. Oder auch l?nger. Und kam Teilweise oben schon mal vor. So wie dieses hier: 'Im folgenden Kapitel sollen die physikalischen Grundlagen erl?utert werden, auf denen die Arbeit basiert'.  Und dann geht es wieder von norne los. Hier kommt jetzt massenhaft Text. Der ist viel. Und kurz. Oder auch l?nger. Und kam Teilweise oben schon mal vor. So wie dieses hier: 'Im folgenden Kapitel sollen die physikalischen Grundlagen erl?utert werden, auf denen die Arbeit basiert'.  Und dann geht es wieder von norne los. Hier kommt jetzt massenhaft Text. Der ist viel. Und kurz. Oder auch l?nger. Und kam Teilweise oben schon mal vor. So wie dieses hier: 'Im folgenden Kapitel sollen die physikalischen Grundlagen erl?utert werden, auf denen die Arbeit basiert'.  Und dann geht es wieder von norne los. Hier kommt jetzt massenhaft Text. Der ist viel. Und kurz. Oder auch l?nger. Und kam Teilweise oben schon mal vor. So wie dieses hier: 'Im folgenden Kapitel sollen die physikalischen Grundlagen erl?utert werden, auf denen die Arbeit basiert'.  Und dann geht es wieder von norne los. Hier kommt jetzt massenhaft Text. Der ist viel. Und kurz. Oder auch l?nger. Und kam Teilweise oben schon mal vor. So wie dieses hier: 'Im folgenden Kapitel sollen die physikalischen Grundlagen erl?utert werden, auf denen die Arbeit basiert'.  Und dann geht es wieder von norne los. Hier kommt jetzt massenhaft Text. Der ist viel. Und kurz. Oder auch l?nger. Und kam Teilweise oben schon mal vor. So wie dieses hier: 'Im folgenden Kapitel sollen die physikalischen Grundlagen erl?utert werden, auf denen die Arbeit basiert'.  Und dann geht es wieder von norne los. 


Hier kommt jetzt massenhaft Text. Der ist viel. Und kurz. Oder auch l?nger. Und kam Teilweise oben schon mal vor. So wie dieses hier: 'Im folgenden Kapitel sollen die physikalischen Grundlagen erl?utert werden, auf denen die Arbeit basiert'.  Und dann geht es wieder von norne los. Hier kommt jetzt massenhaft Text. Der ist viel. Und kurz. Oder auch l?nger. Und kam Teilweise oben schon mal vor. So wie dieses hier: 'Im folgenden Kapitel sollen die physikalischen Grundlagen erl?utert werden, auf denen die Arbeit basiert'.  Und dann geht es wieder von norne los. Hier kommt jetzt massenhaft Text. Der ist viel. Und kurz. Oder auch l?nger. Und kam Teilweise oben schon mal vor. So wie dieses hier: 'Im folgenden Kapitel sollen die physikalischen Grundlagen erl?utert werden, auf denen die Arbeit basiert'.  Und dann geht es wieder von norne los. Hier kommt jetzt massenhaft Text. Der ist viel. Und kurz. Oder auch l?nger. Und kam Teilweise oben schon mal vor. So wie dieses hier: 'Im folgenden Kapitel sollen die physikalischen Grundlagen erl?utert werden, auf denen die Arbeit basiert'.  Und dann geht es wieder von norne los. Hier kommt jetzt massenhaft Text. Der ist viel. Und kurz. Oder auch l?nger. Und kam Teilweise oben schon mal vor. So wie dieses hier: 'Im folgenden Kapitel sollen die physikalischen Grundlagen erl?utert werden, auf denen die Arbeit basiert'.  Und dann geht es wieder von norne los. Hier kommt jetzt massenhaft Text. Der ist viel. Und kurz. Oder auch l?nger. Und kam Teilweise oben schon mal vor. So wie dieses hier: 'Im folgenden Kapitel sollen die physikalischen Grundlagen erl?utert werden, auf denen die Arbeit basiert'.  Und dann geht es wieder von norne los. Hier kommt jetzt massenhaft Text. Der ist viel. Und kurz. Oder auch l?nger. Und kam Teilweise oben schon mal vor. So wie dieses hier: 'Im folgenden Kapitel sollen die physikalischen Grundlagen erl?utert werden, auf denen die Arbeit basiert'.  Und dann geht es wieder von norne los. Hier kommt jetzt massenhaft Text. Der ist viel. Und kurz. Oder auch l?nger. Und kam Teilweise oben schon mal vor. So wie dieses hier: 'Im folgenden Kapitel sollen die physikalischen Grundlagen erl?utert werden, auf denen die Arbeit basiert'.  Und dann geht es wieder von norne los. Hier kommt jetzt massenhaft Text. Der ist viel. Und kurz. Oder auch l?nger. Und kam Teilweise oben schon mal vor. So wie dieses hier: 'Im folgenden Kapitel sollen die physikalischen Grundlagen erl?utert werden, auf denen die Arbeit basiert'.  Und dann geht es wieder von norne los. Hier kommt jetzt massenhaft Text. Der ist viel. Und kurz. Oder auch l?nger. Und kam Teilweise oben schon mal vor. So wie dieses hier: 'Im folgenden Kapitel sollen die physikalischen Grundlagen erl?utert werden, auf denen die Arbeit basiert'.  Und dann geht es wieder von norne los. Hier kommt jetzt massenhaft Text. Der ist viel. Und kurz. Oder auch l?nger. Und kam Teilweise oben schon mal vor. So wie dieses hier: 'Im folgenden Kapitel sollen die physikalischen Grundlagen erl?utert werden, auf denen die Arbeit basiert'.  Und dann geht es wieder von norne los. 
  % Original-Beispielkapitel der Vorlage
%\include{...}

\backmatter
\singlespacing

%\include{Anhang}


% Literaturverzeichnis
% --------------------

% Hier Literatur ergänzen, die zwar ins Verzeichnis aufgenommen werden soll,
% aber im Text nicht zitiert wird.
%\nocite{Ritter83,Horn89}
%\nocite{Jaehne99}
\renewcommand*{\chapterheadstartvskip}{\vspace*{3.3\baselineskip}}% Abstand einstellen
\bibliography{Literature,Literature_H2}

%\include{danksagung}

% Lebenslauf (nur für Doktorarbeiten)
% -----------------------------------
%\include{lebenslauf}

% Selbstständigkeitserklärung (nur für Diplomarbeiten)
% ----------------------------------------------------
\chapter*{}
\section*{Erklärung}\label{chapter:Erklaerung}

Ich versichere, dass ich die Arbeit ohne fremde Hilfe und ohne Benutzung anderer als der angegebenen Quellen angefertigt habe und dass die Arbeit in gleicher oder ähnlicher Form noch keiner anderen Prüfungsbehörde vorgelegen hat und von dieser als Teil einer Prüfungsleistung angenommen wurde. Alle Ausführungen, die wörtlich oder sinngemäß übernommen wurden, sind als solche gekennzeichnet.\newline \newline

\begin{tabularx}{\textwidth}{lXc}
Erlangen, den   &     & \\
\hhline{~~-}
 	  							&     &$\quad$ $\qquad$ Unterschrift Studierende/r $\qquad$$\quad$
\end{tabularx}
\newline
\newline
\newline
\newline
\newline


\chapter*{}
\section*{Erklärung zur Veröffentlichung des Titels einer am Lehrstuhl für Technische Thermodynamik (LTT) durchgeführten Abschlussarbeit}\label{chapter:Erklaerung_Veroeffentlichung}


Name:\rule{0.45\textwidth}{0.5pt} $\qquad$Matrikelnummer:\rule{0.25\textwidth}{0.4pt}  \\
\\
Titel der Arbeit:\rule{0.825\textwidth}{0.4pt} \\
\\
\rule{\textwidth}{0.4pt} \\

Hiermit stimme ich zu, dass der Titel meiner abgeschlossenen \textbf{Master-/ Bachelor-/ Projektarbeit }im Zusammenhang mit meinem Namen
\begin{itemize}
 \item auf den Internetseiten des LTT (http://www.ltt.uni-erlangen.de),
 \item auf gedruckten Aushängen im Lehrstuhlgebäude des LTT sowie
 \item im Newsletter \glqq LTT-aktuell\grqq \, (in gedruckter Form und auf den Internetseiten des LTT per pdf-Download zur Verfügung gestellt) 
\end{itemize}
veröffentlicht werden darf.\\
(nicht Zutreffendes bitte streichen) \newline
\\
\\
\\

\begin{tabularx}{\textwidth}{lXc}
   &     &  \\ \hhline{-~-}
 $\quad$ $\qquad$ Ort, Datum $\qquad$$\quad$	 &     &$\quad$ $\qquad$ Unterschrift Studierende/r $\qquad$$\quad$
\end{tabularx} 	  						



\chapter*{}
\section*{Erklärung zur Nennung des Titels einer extern durchgeführten Abschlussarbeit durch den betreuenden Lehrstuhl für Technische Thermodynamik (LTT)}\label{chapter:Erklaerung_Veroeffentlichung}


Name:\rule{0.45\textwidth}{0.4pt} $\qquad$Matrikelnummer:\rule{0.25\textwidth}{0.4pt}  \\
\\
Titel der Arbeit:\rule{0.825\textwidth}{0.4pt} \\
\\
\rule{\textwidth}{0.4pt} \\
\\
Betreuendes Unternehmen: \rule{0.715\textwidth}{0.4pt} \\
\\
\rule{\textwidth}{0.4pt} \\

Hiermit stimme ich zu, dass der Titel meiner abgeschlossenen \textbf{Master-/ Bachelor-/ Projektarbeit }im Zusammenhang mit meinem Namen und dem des betreuenden Unternehmens
\begin{itemize}
 \item auf den Internetseiten des LTT (http://www.ltt.uni-erlangen.de),
 \item auf gedruckten Aushängen im Lehrstuhlgebäude des LTT sowie
 \item im Newsletter \glqq LTT-aktuell\grqq \, (in gedruckter Form und auf den Internetseiten des LTT per pdf-Download zur Verfügung gestellt) 
\end{itemize}
veröffentlicht werden darf.\\
(nicht Zutreffendes bitte streichen) \newline
\\
\\
\\

\begin{tabularx}{\textwidth}{lXc}
   &     &  \\ \hhline{-~-}
 $\quad$ $\qquad$ Ort, Datum $\qquad$$\quad$	 &     &$\quad$ $\qquad$ Unterschrift Studierende/r $\qquad$$\quad$
\end{tabularx} 	  						
 \\
 \newline
 \\
Zustimmung des betreuenden Unternehmens, vertreten durch: \rule{0.35\textwidth}{0.4pt}
 \newline
  \newline
\begin{tabularx}{\textwidth}{lXc}
   &     &  \\ \hhline{-~-}
 $\quad$ $\qquad$ Ort, Datum $\qquad$$\quad$	 &     &$\quad$ $\qquad$ Unterschrift Unternehmensvertreter/in $\qquad$$\quad$
\end{tabularx} 	


\end{document}


%%% Local Variables:
%%% mode: latex
%%% TeX-master: t
%%% End:
